% Standalone problem document, generated from the adjacent README.md.
% Compile with XeLaTeX. Mathematical content is maintained in Markdown.
\documentclass[11pt,a4paper]{article}
\usepackage[fontsize=10.5pt]{scrextend}
\usepackage[margin=22mm,headheight=15pt,headsep=9mm,footskip=11mm]{geometry}
\usepackage{fontspec}
\setmainfont{texgyrepagella}[Extension=.otf,UprightFont=*-regular,BoldFont=*-bold,ItalicFont=*-italic,BoldItalicFont=*-bolditalic]
\setsansfont{texgyreheros}[Extension=.otf,UprightFont=*-regular,BoldFont=*-bold,ItalicFont=*-italic,BoldItalicFont=*-bolditalic]
\setmonofont{texgyrecursor}[Extension=.otf,UprightFont=*-regular,BoldFont=*-bold,ItalicFont=*-italic,BoldItalicFont=*-bolditalic,Scale=MatchLowercase]
\usepackage{amsmath,amssymb}
\usepackage{unicode-math}
\setmathfont{texgyrepagella-math.otf}
\usepackage{xcolor,microtype,enumitem,needspace,fancyhdr}
\usepackage{bookmark}
\definecolor{ink}{HTML}{17344B}
\definecolor{accent}{HTML}{197D80}
\definecolor{muted}{HTML}{596773}
\hypersetup{colorlinks=true,linkcolor=accent,urlcolor=accent,pdftitle={MI-10 - Chollet's
permanent inequality for Hadamard
products},pdfauthor={Open Problems in Numerical Linear Algebra}}
\urlstyle{same}
\setlength{\parindent}{0pt}
\setlength{\parskip}{5pt plus 1pt minus 1pt}
\linespread{1.04}
\setlength{\emergencystretch}{3em}
\widowpenalty=10000
\clubpenalty=10000
\displaywidowpenalty=10000
\allowdisplaybreaks[1]
\setlist{nosep,leftmargin=1.5em}
\providecommand{\tightlist}{\setlength{\itemsep}{0pt}\setlength{\parskip}{0pt}}
\providecommand{\pandocbounded}[1]{#1}
\pagestyle{fancy}
\fancyhf{}
\fancyhead[L]{\sffamily\footnotesize\color{muted}OPEN PROBLEMS IN NUMERICAL LINEAR ALGEBRA}
\fancyhead[R]{\sffamily\bfseries\footnotesize\color{accent}MI-10}
\fancyfoot[L]{\parbox[b]{0.9\textwidth}{\sffamily\fontsize{8}{10}\selectfont\color{muted}Editorial ratings; a bounded search cannot certify that a problem remains unsolved.\\Literature check: 2026-09-08}}
\fancyfoot[R]{\sffamily\footnotesize\color{muted}\thepage}
\renewcommand{\headrulewidth}{0.3pt}
\renewcommand{\headrule}{\hbox to\headwidth{\color{accent}\leaders\hrule height \headrulewidth\hfill}}
\makeatletter
\renewcommand{\section}{\Needspace{5\baselineskip}\@startsection{section}{1}{0pt}{12pt plus 2pt minus 2pt}{4pt}{\sffamily\bfseries\large\color{ink}}}
\renewcommand{\subsection}{\Needspace{4\baselineskip}\@startsection{subsection}{2}{0pt}{9pt plus 1pt minus 1pt}{3pt}{\sffamily\bfseries\normalsize\color{ink}}}
\makeatother
\setcounter{secnumdepth}{0}
\begin{document}
{\sffamily\small\color{muted}Matrix inequalities and norms\par}
\vspace{3pt}
{\raggedright\hyphenpenalty=10000\exhyphenpenalty=10000\sffamily\bfseries\fontsize{21}{24}\selectfont\color{ink}Chollet's
permanent inequality for Hadamard products\par}
\vspace{7pt}
{\sffamily\small\textbf{Difficulty:} extreme\quad\textbf{Importance:} interesting
to the community\par}
\vspace{4pt}
{\color{accent}\hrule height 0.8pt}
\vspace{6pt}
\textbf{Status:} explicit conjecture; low-order proof claims do not
settle the general case

For an \(n\times n\) matrix \(C=(c_{ij})\), define \[
\operatorname{per}C=\sum_{\sigma\in S_n}\prod_{i=1}^n c_{i,\sigma(i)}.
\] Is it true that, for every integer \(n\ge1\) and every pair of
complex Hermitian positive-semidefinite matrices
\(A,B\in\mathbb C^{n\times n}\), \[
\operatorname{per}(A\circ B)\le
(\operatorname{per}A)(\operatorname{per}B),
\qquad (A\circ B)_{ij}=a_{ij}b_{ij}?
\]

\subsection{Relevance}\label{relevance}

The Hadamard product preserves positive semidefiniteness and corresponds
to tensor products of Gram vectors. The conjecture seeks a computable
structural upper bound for a difficult matrix polynomial under this
basic matrix operation. It concerns matrix inequalities and permanent
computation.

\subsection{References}\label{references}

\begin{itemize}
\tightlist
\item
  J. Chollet, \emph{Is there a permanental analogue to Oppenheim's
  inequality?}, American Mathematical Monthly 89 (1982), 57--58
  (\href{https://doi.org/10.1080/00029890.1982.11995380}{original
  paper}).
\item
  P. Pant and R. Singh, \emph{On Chollet's Permanent Conjecture for
  Graph Laplacians}, arXiv:2604.24192 (2026), abstract and introduction,
  with structured partial results
  (\href{https://arxiv.org/abs/2604.24192}{primary manuscript}).
\item
  Q. Li, Q. Zhao, and W. Luo, \emph{Chollet's Permanent Conjecture
  Through Order Six via Border Induction}, posted 31 August 2026,
  Theorem 1.1
  (\href{https://www.preprints.org/manuscript/202608.2196}{primary
  preprint}).
\item
  I. M. Wanless, \emph{Lieb's permanental dominance conjecture} (2022),
  Conjecture 5 and the implication diagram distinguishing it from
  disproved stronger conjectures
  (\href{https://arxiv.org/pdf/2202.01867}{primary manuscript}).
\end{itemize}

\subsection{Status check --- 2026-09-08}\label{status-check-2026-09-08}

The August 2026 preprint claims the displayed inequality through order
six. That claim has not been independently refereed here and does not
claim the unrestricted result; the unresolved target includes orders
\(n\ge7\). The April 2026 paper addresses structured classes. Searches
for ``Chollet conjecture'', ``permanent'', ``proof'', and 2025--2026
found no general resolution. Disproofs of the stronger Bapat--Sunder and
permanent-on-top conjectures do not refute this statement.
\end{document}
